\documentclass{article}
\usepackage{amsmath, amssymb, amsthm}

\setlength{\parindent}{0pt}

\newtheorem*{claim}{Claim}

\begin{document}

\begin{claim}
    Let $G$ be a connected graph. $G$ has an Eulerian circuit if and only if every vertex of $G$ has even degree.
\end{claim}

\begin{proof}
    ($\Longrightarrow$) Assume $G$ has an Eulerian circuit $C$.

    \medskip

    Consider any arbitrary vertex $v$. Every time the circuit arrives at $v$, it must also leave $v$ along a 
    different edge. Thus the edges incident to $v$ can be paired up as one edge used to enter $v$ and one edge 
    used to leave $v$. Therefore the incident edges occur in pairs, meaning $\deg(v)$ is even. Since $v$ is 
    arbitrary, every vertex of $G$ has even degree.

    \medskip

    ($\Longleftarrow$) Assume every vertex of $G$ has even degree. We use strong induction on the number of edges 
    in $G$.

    \medskip

    \textbf{Base case:} If $G$ has no edges, then since $G$ is connected, it consists of a single vertex. The 
    trivial circuit is an Eulerian circuit.

    \medskip

    \textbf{Inductive step:} Suppose that every connected graph with at most $n$ edges and even degree vertices 
    has an Eulerian circuit. Let $H$ be a connected graph with $n+1$ edges and even degree vertices. 

    \medskip
    
    Since every vertex has even degree and $H$ has edges, every vertex has degree at least $2$. Thus, starting 
    from any vertex and following a path without repeating vertices, since $H$ is finite and every vertex has 
    degree at least $2$, eventually we must encounter a previously visited vertex, which gives a cycle $C$.

    \medskip
    
    Now, remove the edges of $C$ from $H$, leaving a graph $H'$. Let $H_1,\dots,H_k$ be the connected components 
    of $H'$. Removing $C$ decreases the degree of each vertex of $C$ by $2$, so every vertex in each $H_i$ still 
    has even degree. Moreover, each $H_i$ has fewer than $n+1$ edges, so by the inductive hypothesis, each $H_i$ 
    has an Eulerian circuit.

    \medskip

    Since $H$ was connected, each $H_i$ must contain at least one vertex of $C$. At such a common vertex, we can 
    traverse $C$ until reaching that vertex, traverse the Eulerian circuit of $H_i$, and then continue along $C$. 
    Repeating this for each $H_i$, we obtain an Eulerian circuit for $H$.

    \medskip

    Hence $G$ has an Eulerian circuit if and only if every vertex of $G$ has even degree.
\end{proof}

\end{document}
